\subsection{Deterministic Examples}\label{sec:deterministic_scms}

Before giving a formal definition of SCMs, let us revisit the example regarding chocolate consumption and Nobel prizes
and try to come up with quantitative causal models for the data generating process.

For a given country, let us consider two real-valued variables: annual chocolate consumption in kilograms per capita ($C$), and the number of Nobel prize winners per year per capita ($N$). We consider the following simplified hypothetical causal models:
\begin{enumerate}[label=(\roman*)]
  \item $N$ causes $C$ (``Nobel prize winning countries celebrate with massive chocolate consumption''). 
    A tacit but important assumption that underlies this model is that $N$ is not caused by $C$. 
    Assuming an affine relationship, we can model this with the \emph{structural equation}
    \begin{equation}\label{eq:scm_N_causes_C}
    C = \alpha + \beta N
    \end{equation}
    where $\alpha, \beta \in \RN$ are two parameters.
%    For didactical reasons, we use here the ``$\gets$'' symbol to denote that the value of $C$ is \emph{set} to the value of $\alpha + \beta N$, similar to assignment operators in programming languages.\footnote{E.g., in \texttt{R} the assignment operator is \texttt{<-}, in \texttt{Pascal} it is \texttt{:=}, and in \texttt{C} and \texttt{python} it is \texttt{=}.} 
%    In most literature, you will see the notation 
%    $$C = \alpha + \beta N$$ 
%    instead, 
    The convention for such a structural equation is that there must be a single variable on the l.h.s.\ of the equality sign, which indicates that this is the structural equation corresponding to that variable. In the absence of causal cycles,
    the value of the variable on the l.h.s.\ is set to the value of the expression on the r.h.s.\ of the equality sign.
    %(i.e., similarly as in popular programming languages). 
    %We here prefer the ``$\gets$'' notation because it more clearly expresses the asymmetry between cause and effect.
    While the ordinary (``non-structural'') equation $C = \alpha + \beta N$ is usually considered to be equivalent to the equation $N = (C - \alpha) / \beta$ in the sense that their solutions are the same (at least for $\beta \ne 0$), the structural equation $C = \alpha + \beta N$ has an entirely different meaning than the structural equation $N = (C - \alpha) / \beta$ (even for $\beta \ne 0$).
    We will make this difference more precise when we discuss the causal semantics in terms of interventions in the next section.
    For now, we will just note that in the absence of causal cycles, the effect is on the l.h.s.\ of the structural equation, and its direct causes appear on the r.h.s..

    In this model, the variable $C$ is \emph{endogenous}, i.e., its structural equation \eref{eq:scm_N_causes_C} describes how its value is determined by other variables in the model (only $N$ in this case) and model parameters ($\alpha$ and $\beta$ in this case).
    The variable $N$ is \emph{exogenous}, since its value is determined by something external to our model: the model remains silent about what determines the value of $N$.

    The model does make two tacit important assumptions about how the values of $N$ are determined that are not apparent from equation \eref{eq:scm_N_causes_C}. These are:
    \begin{enumerate}[label=\arabic*)]
      \item the exogenous variable $N$ is \emph{not caused by} the endogenous variable $C$, and 
      \item there is \emph{no common cause} of $N$ and $C$. 
    \end{enumerate}
    The first assumption roughly means that if we were to intervene on chocolate consumption $C$ (e.g., by increasing advertizing for chocolate), this would not lead to a change of $N$. The second assumption roughly means that any factor (apart from $N$ and $C$) that causes $N$ does not cause $C$ in any other way than through its effect on $N$. In particular, this means that the parameters $\alpha$ and $\beta$ are assumed to not depend on the value of $N$. In reality, this second assumption is likely to be violated, as the wealth of a country probably causes both $C$ and $N$ directly.

    One concrete way to think about the model (and a very useful analogy) is as a small computer program, for example in the language \texttt{R}:
    \begin{lstlisting}
    N_causes_C <- function(N) {
      C <- alpha + beta * N
      return(C)
    }
    \end{lstlisting}
    The function takes the value of exogenous variable $N$ as input, applies the structural equation \eref{eq:scm_N_causes_C} to calculate the value of endogenous variable $C$, and returns that as its output. The variable $N$ is ``exogenous'' to the function (its value is set externally), while $C$ is ``endogenous'' (its value is set internally).
The computer program analogy also shows that changing the value of $C$ by adapting the code inside the function body does not lead to a change in $N$.
  \item $C$ causes $N$ (``chocolate contains brain enhancing chemicals''). 
    We can model this with the structural equation
    \begin{equation}\label{eq:scm_C_causes_N}
    N = \gamma + \delta C
    \end{equation}
    with two parameters $\gamma,\delta \in \RN$.
    For this model, we again distinguish exogenous and endogenous variables, 
    but now $C$ is considered the exogenous variable, and $N$ the endogenous variable.
    A computer program representation of this model in \texttt{R} is:
    \begin{lstlisting}
    C_causes_N <- function(C) {
      N <- gamma + delta * C
      return(N)
    }
    \end{lstlisting}
    The tacit assumptions are that $N$ does not cause $C$, and that $N$ and $C$ are not confounded.
  \item $C$ and $N$ have a common cause, wealth $W$ (``The wealth of a country determines both how much money goes to science and also how much people can spend on chocolate'').
    We could model this using two structural equations:
    \begin{align}\label{eq:scm_W_causes_C_and_N}
      C &= \alpha + \beta W \\
      N &= \gamma + \delta W
    \end{align} 
    In this case, we treat the variables $C$ and $N$ both as endogenous, and we treat $W$ as exogenous.
    The following tacit assumptions are made:
    \begin{enumerate}[label=\arabic*)]
      \item $W$ is not caused by $C$, nor by $N$;
      \item none of the variables is confounded, i.e., there is no other variable that is a common cause of any pair of $\{N,C,W\}$, or even the triple $(N,C,W)$.
    \end{enumerate}
    Another way to formulate the second assumption is that the values of $W$, the parameters $(\alpha,\beta)$ and the parameters $(\gamma,\delta)$ are all determined ``independently''.

    A computer program representation of this model in \texttt{R} is:
    \begin{lstlisting}
    W_causes_C_and_N <- function(W) {
      C <- alpha + beta * W
      N <- gamma + delta * W
      return(c(C,N))
    }
    \end{lstlisting}
    It takes as input the value of the exogenous variable $W$, applies the structural equations to calculate the values of the endogenous variables $C$ and $N$ accordingly, and returns a tuple with both values as output.
\end{enumerate}
Even with only two (or three) variables, many different causal models can be made, and the above three illustrate just a small subset of those. Indeed, we could for example combine the third model with a direct causal effect of $C$ on $N$ (or the other way around), and we could assume more complicated confounding relationships involving latent confounders between the variables.
In case you wonder how to model the other causal hypotheses (cyclic relationship, selection bias, and functional constraints):
we will postpone discussion of these since they are inherently more complicated.

\subsection{Hard Interventions}

So far, we have been rather informal about the causal meaning of the models (i), (ii) and (iii), which provide different causal explanations for the observed correlation between $C$ and $N$. 
We will now make this more precise in terms of hypothetical \emph{interventions} that we could perform on the ``system'' (a country) consisting of the two variables $C$ and $N$.

Let us start again with discussing the first model ($N$ causes $C$), with structural equation \eref{eq:scm_N_causes_C}.
We consider two interventions:
\begin{itemize}
  \item What would happen if Putin forced the Russians to eat $c$ kilograms of chocolate every year?
    We can model this \emph{intervention on (the structural equation for) $C$} by modifying structural equation \eref{eq:scm_N_causes_C} into the following intervened structural equation:
    \begin{equation}\label{eq:scm_N_causes_C_do_C}
    C = c.
    \end{equation}
    This intervention would obviously lead to a change of the value of $C$ (for most choices of $c$, i.e., except if $c$ happens to equal $\alpha + \beta N$). It would, however, have no effect on the number of Russian Nobel prize winners: the value of $N$ would remain unaltered, since exogenous variables (in this model, $N$) are assumed not to be caused by endogenous ones (for this model, $C$).
    
    Following Pearl, we will denote this intervention as $\doit(C=c)$. It is often called a \emph{hard} (or ``\emph{surgical}'', or ``\emph{atomic}'', or ``\emph{perfect}'') intervention because it completely overrides the default causal mechanism that normally determines the value of $C$ expressed by structural equation \eref{eq:scm_N_causes_C}.

    The computer program equivalent of this intervention is to replace the line of code \texttt{C <- alpha + beta * N} with the assignment \texttt{C <- c}. If we consider $c$ to be a (constant) parameter, we get the following ``intervened'' program:
    \begin{lstlisting}
    N_causes_C.do_c <- function(N) {
      C <- c
      return(C)
    }
    \end{lstlisting}
  \item
    What would happen if we somehow enforced that the Russians win a certain number of Nobel prizes (e.g., by manipulating the Nobel prize committee members)?
    In this case, the causal mechanism determining the chocolate consumption modeled by structural equation \eref{eq:scm_N_causes_C} would still be in place. Therefore, changing the exogenous value $N$ would generically (i.e., if $\beta \ne 0$) lead to a change of $C$. In other words, if before the intervention the chocolate consumption is $c = \alpha + \beta n$, and we intervene to set $N=n'$, then after the intervention the value of $C$ will become $c' = \alpha + \beta n'$. This hard intervention is denoted as $\doit(N=n')$.
    It leaves the structural equation \eref{eq:scm_N_causes_C} invariant. However, to
    reflect that the value of $N$ is now set to $n'$, we will add the structural equation
    $$N = n'$$
    to the model, and from now on will treat $N$ as endogenous to reflect that its value has been set.
    Thus, the intervened model has structural equations:
    \begin{align*}
      N & = n' \\
      C & = \alpha + \beta N
    \end{align*}
    and treats both $C$ and $N$ as endogenous variables.
% We actually have two modeling options here. The first is to add a structural equation $N \to n'$ and from now on, treat $N$ as endogenous (with its value hardcoded to $n'$ as a result of the intervention). This corresponds with the intervention $\doit(N=n)$. The second is to still treat $N$ as exogenous, and simply keep using the original structural causal model, which does not specify the value of $n'$. The latter corresponds with the intervention (family) $\doit(N)$.
\end{itemize}
We summarize this discussion in terms of a larger computer program example (in Figure~\ref{fig:scm_N_causes_C}) that illustrates what this causal model predicts under different interventions.

\begin{figure}\centering
\begin{lstlisting}
### observational model
alpha <- 2.0
beta <- 0.5
N_causes_C <- function(N) {
  C <- alpha + beta * N
  return(C)
}
# query model for the Dutch (n=11):
N_causes_C(11)

### hard intervention: do (C=c)
c <- 12 # observed value for Switzerland
N_causes_C.do_c <- function(N) {
  C <- c
  return(C)
}
# query model for the Dutch (n=11):
N_causes_C.do_c(11)

### hard intervention: do (N=n)
n <- 32 # observed value for Switzerland
N_causes_C.do_n <- function() {
  N <- n
  C <- alpha + beta * N
  return(c(C,N))
}
# query model:
N_causes_C.do_n()
\end{lstlisting}
\caption{Computer program illustrating how to simulate from the causal model ``$N$ causes $C$'' under various interventions.\label{fig:scm_N_causes_C}}
\end{figure}

\begin{Exc}
  Write an analogous \texttt{R} program for the second model, $C$ causes $N$. Remember that it
  treats $C$ as exogenous, $N$ as endogenous, and has structural equation \eref{eq:scm_C_causes_N}.
  Simulate from the model in the observational setting, after the hard intervention
  $\doit(C=c)$ and after the hard intervention $\doit(N=n)$.
\end{Exc}

For what is presumably the most realistic model ($C$ and $N$ are caused by $W$) with structural equations
\eref{eq:scm_W_causes_C_and_N} and where $W$ is treated as exogenous, and $C$ and $N$ as endogenous, we discuss three hard interventions:
\begin{itemize}
  \item $\doit(C=c)$. This changes the structural equations into:
    \begin{align}\label{eq:scm_W_causes_C_and_N_do_C}
      C &= c \\
      N &= \gamma + \delta W
    \end{align} 
    Note that while the structural equation for $C$ is changed by the intervention, the structural equation for $N$ is not (i.e., it remains invariant). This reflects the ``perfect'' character of the intervention: it only changes the structural equation(s) of the intervention target(s) (only $C$ in this case), without any side effects on other structural equations.
  \item $\doit(N=n)$. Similarly, this only changes the structural equation for $N$, without modifying the structural equation for $C$, and yields the following intervened structural equations:
    \begin{align}\label{eq:scm_W_causes_C_and_N_do_N}
      C &= \alpha + \beta W \\
      N &= n
    \end{align} 
  \item $\doit(W=w)$. Since $W$ is exogenous, the structural equations for $C$ and $N$ remain the same as in the ``observational'' (pre-interventional) setting, i.e., as in \eref{eq:scm_W_causes_C_and_N}, and we add a structural equation for $W$:
    \begin{align*}
      C &= \alpha + \beta W \\
      N &= \gamma + \delta W \\
      W &= w
    \end{align*}
  In this intervened model, all three variables are treated as endogenous variables.
\end{itemize}

%We finish by giving a deterministic version of Definition~\ref{def-causality-real}:
%\begin{Def}[Causal effect---deterministic version]
%    \label{def:causality-real-deterministic}
%    We say that endogenous variable $X$ has \emph{causal effect} on another endogenous variable $Y$ if
%    under the hard intervention $\doit(X=x)$, the value of $Y$ depends on $x$. In other words,
%    if there exist two values $x,x'$ such that $Y_x \ne Y_{x'}$.
%\end{Def}

\begin{Exc}
  For each of the three hard interventions considered above: what variable values are changed due to the intervention (comparing with the pre-interventional ``observational'' model)? 
  How do these correspond to the intuitive causal effects between the three variables?
\end{Exc}

\subsection{Soft Interventions}

A more general class of interventions are \emph{soft interventions}. These modify 
the expression on the r.h.s.\ of a structural equation, but without adding new variables
to that expression. Special cases are hard interventions on endogenous variables.

As an example, consider the model with the confounder $W$ for $C$ and $N$, with 
structural equations:
    \begin{align*}
      C &= \alpha + \beta W \\
      N &= \gamma + \delta W.
    \end{align*}
An example of a soft intervention on $C$ is a change in the parameter $\beta$, replacing it
with $\beta'$ (perhaps to reflect a change in the chocolate price):
    \begin{align*}
      C &= \alpha + \beta' W \\
      N &= \gamma + \delta W.
    \end{align*}
Another example would be to change the functional form of the structural equation:
    \begin{align*}
      C &= \beta\gamma + \alpha W^3 \\
      N &= \gamma + \delta W.
    \end{align*}
In contrast to a hard intervention on $C$, which would remove all occurrences of 
endogenous and exogenous variables at the r.h.s.\ of the structural equation for $C$,
this soft intervention only changes the functional / parametric form of the equation,
but its value may still depend on the variables that appeared in the original 
structural equation.

In general, a soft intervention on $C$ will result in the following structural equations:
    \begin{align*}
      C &= f_{\theta}(W) \\
      N &= \gamma + \delta W
    \end{align*}
for some function $f_{\theta} : \RN \to \RN$ parameterized by some parameter $\theta \in \Theta$.
Similarly, a soft intervention on $N$ will result in
    \begin{align*}
      C &= \alpha + \beta W \\
      N &= g_{\theta}(W)
    \end{align*}
for some function $g_{\theta} : \RN \to \RN$ parameterized by some parameter $\theta \in \Theta$.

\subsection{Intervention Variables}\label{sec:intervention_variables}

It can be very convenient to introduce explicit variables that indicate whether or how
interventions have been performed. These are often referred to as \emph{intervention variables}
(and are special cases of so-called \emph{context variables}, \emph{environment variables} and \emph{domain indicators}).

Often, intervention variables can be considered as exogenous. For instance, if their values (i.e., which intervention
is performed, and how) are determined earlier in time than the values of the endogenous variables, 
which means that the endogenous variables cannot cause the intervention variables (except, perhaps, if
time travel becomes possible). An
example is a scientific experiment in which the experimenter prepares the system in some 
experimental condition before performing measurements, which implies that the experimental condition 
cannot be influenced by the measured state of the system. 
A concrete instance of such an intervention variable is 
the coin flip that determines whether a subject enters the treatment or control group in a
randomized controlled trial. 
However, if a doctor prescribes a certain treatment based on 
the state of the patient determined through a diagnosis, then the treatment variable is 
influenced by the patient's state and should therefore not be considered to be exogenous with
respect to the patient's state.

Let us discuss the randomized controlled trial example as described in Principle~\ref{prn:rct} in more detail.
\begin{Eg}[Binary intervention variable]
Let $X$ be the real-valued variable that describes treatment (e.g., the prescribed dose of a drug), and $Y$ be the real-valued variable that quantitatively describes the outcome.
We consider $X$ and $Y$ as endogenous variables.
Introduce the binary intervention variable $C$ that indicates whether the subject enters the test group ($C=1$) or the control group ($C=0$). 
In practice, one uses a coin or, nowadays, a random number generator, to determine the value of $C$. 
Due to the experimental setup of the RCT, we can therefore assume that both $Y$ and $X$ do not cause $C$. 
Thus we can consider $C$ as an exogenous variable, and can write down the following structural equations to model the RCT:
  \begin{align*}
    X &= \begin{cases}
      x_0 & \text{if $C=0$,} \\
      x_1 & \text{if $C=1$} 
    \end{cases} \\
    Y &= \alpha + \beta X.
  \end{align*}
The parameter $\beta$ is often referred to as \emph{the causal effect of $X$ on $Y$}, in the sense
of being the effect that a unit change of $X$ has on $Y$.
In addition to the assumption that $C$ is not caused by $X$ or $Y$, an important assumption is
also that there is no confounding between $C$ and $Y$. This boils down to assuming that the outcome
of the coin flip is not caused by anything else that also causes the outcome $Y$.
\end{Eg}

%Alternatively, we could consider $X$ itself to be the intervention variable; then there is no
%need to introduce the variable $C$, we would consider $X$ as exogenous and $Y$ as endogenous,
%and the model would only consist of the structural equation for $Y$:
%    $$Y \gets \alpha + \beta X.$$
There is no need for an intervention variable to be binary or discrete.
An example of a continuous intervention variable is the dose of the drug used for treatment.
\begin{Eg}[Continuous intervention variable]
The dose of a drug used for treatment often quantitatively affects the outcome. 
This dependence is visualized in dose-response curves that link a quantitative response to the dose of the drug that was administered. 
The following \emph{Hill model} is often used:
$$Y = \frac{Y_{max}}{1 + \left(\frac{D}{D_{50}}\right)^{-n}}$$
where $Y$ is the response, $D$ the dose, and $Y_{max}$, $D_{50}$ and $n$ are parameters
(the maximum response, the potency, and the Hill coefficient, respectively).
We can consider $D$ to be a continuous (non-negative) intervention variable, that not only indicates whether an intervention was performed ($D > 0$?), but also how the intervention was performed (the administered dose).
It is again obvious from the setting that the response does not causally influence the dose.
However, if there is no explicit and careful randomization, one cannot always assume that there is no common cause of dose and response. 
For example, if physicians decide to not treat patients with terminal cancer because of the strong side effects of the treatment, then the stage of the cancer is a common cause of dose and response.
\end{Eg}
%\footnote{However, 
%one has to be careful about implicit temporal aspects: if the response $Y$ refers
%to the response of a previous treatment, and $D$ is the dose of a next treatment, then it
%could be very well the case that $Y$ causes $D$, for example, if a doctor is in the process
%of adjusting the dose to obtain an optimal response.}


\subsection{Modeling Uncertainty}

So far, we have not introduced any randomness into our causal models. This is
often an important step when modeling a system. While this is a \emph{sine qua
non} for purely probabilistic/statistical modeling, it is optional---yet
common---in causal modeling. 

The main reason for introducing stochasticity into a model is to deal with
missing information about the system state. For example, when
rolling a die, if one would have access to the exact initial state (position, momentum,
angular momentum) of the die once released, one could calculate the outcome of the 
roll as a (complicated) deterministic function of the initial state. In
practice, it is neither feasible to measure or control the initial state of the die
after it is released from the hand, nor to calculate this deterministic function. 
Therefore, instead of attempting to model this
meticulously, one can deal with any remaining uncertainty by modeling its probability distribution
(which could either describe the subjective beliefs about the value of a variable
in a Bayesian interpretation, or, in a frequentist interpretation, a description of 
its statistics given an infinite ensemble consisting of copies of the system). 

In causal modeling, we often have to 
incorporate stochasticity in this way. Therefore, in a structural causal model,
we will specify certain probability distributions on subsets of the \emph{exogenous}
variables to explicitly represent ``noise'' that may affect the values of the
endogenous variables. This leads us to distinguish three types of variables:
endogenous variables, exogenous input variables, and exogenous random variables.
The only difference between the exogenous input variables and the exogenous
random variables is that the model specifies the probability distributions of the
latter, while not making any assumptions on the values assumed by the former.
Often, the exogenous random variables are introduced to model ``noise'' and 
are considered latent; hence they are sometimes referred to as ``noise'' or
``disturbance'' variables.

As an example, we will again revisit the chocolate consumption and Nobel price
winners example. Consider the causal model $C$ causes $N$, based on the hypothesis
that chocolate consumption enhances cognitive abilities. When looking at the data
in Figure~\ref{fig:choco-nobel}, it is obvious that an exact affine relationship
of the form $N = \gamma + \delta C$ cannot be very accurate, as for example
The Netherlands has the same chocolate consumption as Australia, but obtained 
considerably more Nobel prizes per capita. We can hypothesize that 
other factors determine the value of $N$, apart from $C$. Rather than attempting to meticulously
model all of these factors, we will introduce just a single ``effective'' or 
``aggregated'' exogenous random
variable $E$ that expresses their residual influence on $N$. We thus refine our
structural equation for $N$ to:
  \begin{equation}\label{eq:se_C_causes_N_noisy}
  N = \gamma + \delta C + E
  \end{equation}
and in addition we assume that $E$ has a certain probability distribution $P^E$.
For concreteness, we could assume that $E \sim \Ncal(0,\sigma^2)$, i.e., that $E$
has a Gaussian distribution with zero mean and variance $\sigma^2$, where we 
introduced a new parameter $\sigma^2$ to the model. Summarizing, the stochastic
structural causal model then becomes:
  \begin{align*}\label{eq:scm_C_causes_N_noisy}
      E \sim {} & \Ncal(0,\sigma^2) \\
      N = {} & \gamma + \delta C + E,
  \end{align*}
with $C$ an exogenous input variable, $E$ an exogenous random variable,
and $N$ an endogenous variable, and parameters $\gamma,\delta,\sigma^2$.

The model can again be implemented as a computer program in \texttt{R}:
    \begin{lstlisting}
    C_causes_N_stoch <- function(C) {
      E <- rnorm(1,mean = 0,sd = sigma)
      N <- gamma + delta * C + E
      return(N)
    }
    \end{lstlisting}
This function takes as input the value of $C$, then samples an (independent) random 
value for $E$, and uses that to calculate the value for $N$ that is returned.
Here, $N$ is the endogenous variable, $C$ is an exogenous input variable, and $E$ is an exogenous random variable.

Note that the analogy is not entirely accurate, as the structural causal model
prescribes the probability distribution of $E$, while the computer program samples 
a value from this distribution. The reason is that a language like \texttt{R} does
not allow us to express probability distributions; for that, we would need a 
probabilistic programming language, but those are less well known and therefore less
useful for our didactic purpose. Henceforth we will think of computer implementations
of a structural causal model as programs that \emph{sample} from the model, and 
whose elements correspond directly with those of the model.
